\SourceLeaf{10}
\[
\mathrm{d}\Sigma=-\mathrm{d}s'.\cos QX'=+\mathrm{d}s''.\cos QX''=-\mathrm{d}s'''.\cos QX'''\text{ etc.}
\]
and therefore, since the number of parts is even,
\[
\mathrm{d}s'.\cos QX'+\mathrm{d}s''.\cos QX''+\mathrm{d}s'''.\cos QX'''+\text{ etc.}=0
\]
\par
By treating all the remaining elements $\mathrm{d}\Sigma$ in the same way and summing, we obtain
\par\medskip\Needspace{4\baselineskip}
\begin{center}THEOREM ONE.\end{center}
\par
\textit{The integral $\int\mathrm{d}s\cos QX$ taken over the whole surface of the body is $=0$.}
\par
More generally, the same method shows that the integral
\[
\int(T\cos QX+U\cos QY+V\cos QZ)\mathrm{d}s
\]
vanishes if $T,U,V$ denote rational functions of $y,z$ alone, of $x,z$ alone, and of $x,y$ alone, respectively.

\Article{4}
\par
Since the volumes of the portions of the cylinder from our plane to the respective points $P',P'',P'''$ and so on are $=\mathrm{d}\Sigma.(x'-\alpha)$, $\mathrm{d}\Sigma.(x''-\alpha)$, $\mathrm{d}\Sigma.(x'''-\alpha)$ etc., the part of the volume of the body that lies inside the cylinder will be
\[
\begin{aligned}
&=-x'\mathrm{d}\Sigma+x''\mathrm{d}\Sigma-x'''\mathrm{d}\Sigma+\text{ etc.}\\
&=\mathrm{d}s'.x'\cos QX'+\mathrm{d}s''.x''\cos QX''+\mathrm{d}s'''.x'''\cos QX'''+\text{ etc.}
\end{aligned}
\]
and summing over all $\mathrm{d}\Sigma$ we obtain
\par\medskip\Needspace{4\baselineskip}
\begin{center}THEOREM TWO.\end{center}
\par
\textit{The entire volume of the body is expressed by the integral $\int\mathrm{d}s.x\cos QX$ taken over the whole surface.}
\par
Evidently the same volume can also be expressed by $\int\mathrm{d}s.y\cos QY$ or by $\int\mathrm{d}s.z\cos QZ$.

\Article{5}
\par
First imagine the entire cylinder filled with matter of uniform density, and let us see how much attraction each of its elements exerts on point $M$\SourceLeaf{11}. Divide the cylinder into elementary cylinders by planes infinitely close to one another and parallel to the base. One such elementary cylinder, at the point whose coordinates are $\xi,\eta,\zeta$, can be expressed by $\mathrm{d}\Sigma.\mathrm{d}\xi$. Its distance from point $M$ will be
\[
=\BareRoot{((a-\xi)^2+(b-\eta)^2+(c-\zeta)^2)}=\rho
\]
Hence its attraction on point $M$ can be expressed by $\mathrm{d}\Sigma.\mathrm{d}\xi.f\rho$, where the function $f\rho$ denotes the law of attraction. Since throughout the cylinder only $\xi$ is to be regarded as variable, we have $\rho\mathrm{d}\rho=-(a-\xi)\mathrm{d}\xi$, and consequently the attraction of the element is $=-\frac{\rho f\rho.\mathrm{d}\rho.\mathrm{d}\Sigma}{(a-\xi)}$. Resolving this attraction into three components parallel and opposite to the coordinate axes $x,y,z$, we find that the first is $=-f\rho.\mathrm{d}\rho.\mathrm{d}\Sigma$. Thus, denoting the integral $\int f\rho.\mathrm{d}\rho$ by $F\rho$, the attraction of the cylinder from the base $\mathrm{d}\Sigma$ up to the point whose first coordinate is $=\xi$, exerted on point $M$ along the $x$-coordinate axis, will be $=-(F\rho-\mathrm{Const.})\mathrm{d}\Sigma=-(F\rho-FR)\mathrm{d}\Sigma$, if $R$ is taken to denote the distance of the base $\mathrm{d}\Sigma$ from point $M$. It follows that this same component of the attraction from all the parts of the body lying inside the cylinder is
\[
\begin{aligned}
&=(Fr'-Fr''+Fr'''-\text{ etc.})\mathrm{d}\Sigma\\
&=-Fr'.\mathrm{d}s'.\cos QX'-Fr''.\mathrm{d}s''.\cos QX''-Fr'''.\mathrm{d}s'''.\cos QX'''-\text{ etc.}
\end{aligned}
\]
\par
Extending this reasoning to all elements $\mathrm{d}\Sigma$, we obtain
\par\medskip\Needspace{4\baselineskip}
\begin{center}THEOREM THREE.\end{center}
\par
\textit{The attraction that the body exerts on point $M$ in the direction parallel and opposite to the $x$-coordinate axis is expressed by the integral $-\int Fr.\mathrm{d}s.\cos QX$ taken over the whole surface.}
\par
{\rightskip=0pt plus 3em Evidently, in exactly the same way, the attraction along the other two principal directions is expressed by the integrals $-\int Fr.\mathrm{d}s.\cos QY$, $-\int Fr.\mathrm{d}s.\cos QZ$.\par}

\Article{6}
\par
We will now approach the matter by another route. Imagine a spherical surface of radius $=1$ drawn about the center $M$ and divided into infinitely small elements. Let $\Pi$ be a point on this surface belonging to the small patch $\mathrm{d}\Sigma$ on it. Draw the radius $M\Pi$ and, if necessary, extend it indefinitely beyond the surface of the sphere. Let $P',P'',P''$ etc. be the points at which this radius successively intersects the surface of our body, excluding point $M$ itself if it happens to lie on the surface. The number of these
